\documentclass[aspectratio=169,12pt]{beamer}

\usepackage[T1]{fontenc}
\usepackage[utf8]{inputenc}
\usepackage[english]{babel}
\usepackage{lmodern}
\usepackage{booktabs}
\usepackage{tabularx}

\newcommand{\seriesconn}{\mathbin{\oplus}}
\newcommand{\Lel}[1]{\mathrm{L}\!\left(#1\right)}
\newcommand{\Wel}[1]{\mathrm{W}\!\left(#1\right)}
\newcommand{\Gel}[2]{\mathrm{G}\!\left(#1,#2\right)}
\newcommand{\CPEel}[2]{\mathrm{CPE}\!\left(#1,#2\right)}


\definecolor{ink}{HTML}{19324A}
\definecolor{accent}{HTML}{176E84}
\setbeamercolor{normal text}{fg=ink,bg=white}
\setbeamercolor{frametitle}{fg=ink,bg=white}
\setbeamercolor{structure}{fg=accent}
\setbeamercolor{title}{fg=ink}
\setbeamerfont{frametitle}{size=\LARGE,series=\bfseries}
\setbeamerfont{normal text}{size=\large}
\setbeamerfont{title}{size=\Huge,series=\bfseries}
\setbeamertemplate{navigation symbols}{}
\setbeamertemplate{footline}[frame number]
\setbeamertemplate{itemize items}[circle]

\title{Discovering circuits from impedance data}
\subtitle{SOFCs, equivalent circuits, and formal grammars}
\author{Mitja Mandić}
\institute{IJS E8}
\date{}

\begin{document}

\begin{frame}[plain]
  \titlepage
\end{frame}

\begin{frame}
  {Agenda}
  \tableofcontents
%  \begin{itemize}
%    \item What is a (solid oxide) fuel cell?
%    \item Intro to electrochemical impedance spectroscopy
%    \item Connection to formal grammars
%    \item Our implementation and results
%  \end{itemize}
\end{frame}

\section{What is a (solid oxide) fuel cell?}


\begin{frame}{What is a solid oxide fuel cell?}
  \begin{itemize}
    \setlength{\itemsep}{0.95em}
    \item Fuel cell is an ``electricity factory''
    \item A SOFC converts fuel and oxygen into electricity and heat.
    \item Its hot ceramic electrolyte carries oxygen ions from the cathode to the anode.
    \item At the anode, the ions react with fuel. Released electrons flow through an external circuit. 
  \end{itemize}
\end{frame}

\begin{frame}{What is a solid oxide fuel cell?}
  %include sofc graphic
\end{frame}

\begin{frame}{How is a SOFC different from a battery?}
  \renewcommand{\arraystretch}{1.55}
  \begin{tabularx}{\textwidth}{@{}>{\raggedright\arraybackslash}p{0.19\textwidth}>{\raggedright\arraybackslash}X>{\raggedright\arraybackslash}X@{}}
    \toprule
    \textbf{Aspect} & \textbf{SOFC} & \textbf{Battery} \\
    \midrule
    Fuel supply
      & External fuel and air
      & Reactants stored inside \\
    Runtime
      & Runs while fuel flows
      & Runs until stored energy is used \\
    Continue use
      & Add more fuel
      & Recharge or replace \\
    Operation
      & High temperature, slower startup
      & Usually near room temperature, quick startup \\
    \bottomrule
  \end{tabularx}
\end{frame}
\section{Intro to electrochemical impedance spectroscopy}
\begin{frame}{Impedance and EIS data}
  \begin{itemize}
    \setlength{\itemsep}{0.95em}
    \item \textbf{E}lectrochemical \textbf{I}mpedance \textbf{S}pectroscopy
    applies a small oscillating voltage and measures the current response across frequencies.
    \item The resulting impedance spectrum describes how the cell responds at different frequencies.
    \item High frequencies (10 kHz) show resistive behaviour, lower (1 mHz) capture charge transfer and ion diffusion.
    \item Non-invasive technique, the idea is to capture the data from a cell as is (without changing the system).
    %\item An equivalent circuit combines simple elements to reproduce the spectrum.
  \end{itemize}
\end{frame}

\begin{frame}{Physical vs. equivalent circuit modeling (ECM)}
  \onslide<1->{
  Option 1: Model the cell as a physical system
  \begin{itemize}
    \item Try to capture the processes inside the cell 
    in a system of equations
    \item Requires detailed knowledge of the chemical process
    \item Very complex, lots of variables
  \end{itemize}}
  \vspace{0.5cm}
  \onslide<2>{
    Option 2: Take the data and fit a ``simple'' explanation
    \begin{itemize}
      \item \textit{Data science-y} approach
      \item Fit a model we understand and draw conclusions from it
      \item Chosen technique: equivalent circuit modeling
    \end{itemize}
  }
\end{frame}


\begin{frame}{ECM building blocks}
  \renewcommand{\arraystretch}{1.75}
  \begin{tabularx}{\textwidth}{@{}>{\raggedright\arraybackslash}p{0.19\textwidth}>{\centering}X>{\raggedright\arraybackslash}X@{}}
    \toprule
    \textbf{Element} & \textbf{Impedance formula} & \textbf{Meaning in simple terms} \\
    \midrule
    \textbf{R}\enspace Resistor
      & $Z_R = R$ & Opposes current. Its impedance does not change with frequency. \\
    \textbf{CPE}\enspace Constant phase element
      & $Z_{CPE} = \frac{1}{(j\omega)^{\alpha}Q}$ & Imperfect capacitor and non-ideal charge storage. \\
    \textbf{L}\enspace Inductor
      & $Z_L = j\omega L$ & Inductive behavior at high frequency. \\
    \textbf{W}\enspace Warburg element
      & $Z_W = \frac{\sigma}{\sqrt{jw}}$ & Models diffusion and limits from slow mass or ion transport. \\
    \bottomrule
  \end{tabularx}
\end{frame}
\section{Connection to formal grammars}
\begin{frame}
  {Connection to formal grammars}
\begin{itemize}
  \item Our idea was to create a formal grammar that would create all possible
  circuits
  \item We specify the rules how elements may connect and define how complex 
  circuits do we consider
  \item In the frequency domain impedance equations follow immediately; for
  elements in series we sum their impedances, and in parallel we sum their
  admittances (reciprocal impedance values) 
\end{itemize}

\end{frame}


\begin{frame}{Circuit grammar}
  \begin{center}
    \renewcommand{\arraystretch}{1.45}
    \begin{tabular}{@{}rcl@{}}
      \texttt{Circuit} & \texttt{::=} & \texttt{Rs} \\
                       &               & \texttt{or Rs + Network} \\[0.5em]
      \texttt{Network} & \texttt{::=} & \texttt{Element} \\
                       &               & \texttt{or (Network + Network)} \\
                       &               & \texttt{or (Network || Network)} \\[0.5em]
      \texttt{Element} & \texttt{::=} & \texttt{R, CPE, L, or W} \\
    \end{tabular}
  \end{center}

  \vspace{0.7em}
  \begin{itemize}
    \setlength{\itemsep}{0.5em}
    \item \texttt{+} connects networks in series and \texttt{||} connects them in parallel.
    \item This grammar permits all combinations of elements, 
    with a resistor \texttt{Rs} connected in series
  \end{itemize}
\end{frame}

\section{Our implementation and results}

\begin{frame}{Circuit generation process}
  \footnotesize
  \renewcommand{\arraystretch}{1.10}
  \begin{tabularx}{\textwidth}{@{}r>{\ttfamily\raggedright\arraybackslash}X@{}}
     1 & INPUT maximum element count $N$, primitives or blocks $P$ \\
     2 & $B[1\ldots N-1] \gets$ empty buckets, $C \gets \{\mathrm{Rs}\}$ \\
     3 & for each $p \in P$ \\
     4 & \hspace*{1em}$q \gets \mathrm{Normalize}(p)$, then add $q$ to $B[|q|]$ when $|q| < N$ \\
     5 & for $n = 1\ldots N-1$ \\
     6 & \hspace*{1em}for each split $n = l + r$ with $l \le r$ \\
     7 & \hspace*{2em}for each $(u,v) \in B[l] \times B[r]$ and $o \in \{\mathrm{series},\mathrm{parallel}\}$ \\
     8 & \hspace*{3em}$w \gets \mathrm{Normalize}(o(u,v))$ \\
     9 & \hspace*{3em}if $|w| = n$, add $w$ to $B[n]$ \\
    10 & \hspace*{1em}for each $w \in B[n]$, add $\mathrm{Normalize}(\mathrm{series}(\mathrm{Rs},w))$ to $C$ \\
    11 & RETURN $C$ \\
  \end{tabularx}
\end{frame}

\begin{frame}{Parameter fitting}
  \begin{itemize}
    \setlength{\itemsep}{0.65em}
    \item We fit a predicted impedance $\widehat{Z}(f;\boldsymbol{\theta})$ for each circuit 
    \item Residuals use separate component-relative scales so that no frequency range
      dominates only because its impedance values are larger.
    \item Each circuit uses six initial parameter sets: one fixed start that is the same for all circuits
    and five stratified random starts. We keep the fit with the lowest optimizer cost.
    \item Bounded SciPy \texttt{least\_squares} fits the parameters with the TRF method.
      CPE exponent $\alpha$ remains linear while the other parameters use log space.
    \item The \texttt{soft\_l1} loss reduces the influence of large residuals.
      User defined parameter \texttt{f\_scale} sets when this downweighting begins.
      Best performing value we tested was \texttt{f\_scale}$=10^{-3}$
  \end{itemize}
\end{frame}

\begin{frame}{Candidate screening and refinement}
  \begin{enumerate}
    \setlength{\itemsep}{0.8em}
    \item Fit every generated circuit with the same multi-start procedure and a limited
      first-round evaluation budget.
    \item Rank the first-round fits by optimizer cost and retain circuits within 20\%
      of the best value.
    \item Restart each retained circuit from its fitted parameters, using tighter
      tolerances and a much larger evaluation budget.
    \item Rank the refined circuits with $L_{0.5}$-BIC, which rewards a good fit while
      penalizing additional parameters.
  \end{enumerate}
\end{frame}

\begin{frame}{Simulation study}
Data generated between 100 kHz and 0.1 Hz from following circuits with added noise:
\[
\begin{aligned}
\textnormal{SOFC circuit}
    &= R_s
    \seriesconn
    \left[
        R_1
        \parallel
        \CPEel{Q_2}{\alpha_2}
    \right]
    \seriesconn
    \left[
        R_3
        \parallel
        \CPEel{Q_4}{\alpha_4}
    \right]
    ,
    \\[1.5ex]
\textnormal{Battery circuit}
    &= R_s
    \seriesconn
    \left[
        R_1
        \parallel
        \CPEel{Q_2}{\alpha_2}
    \right]
    \seriesconn
    W(\sigma_3).
\end{aligned}
\]
\onslide<2->{
Battery circuit was correctly identified, however SOFC was not
\[
\textnormal{Selected SOFC circuit}
    = R_s
    \seriesconn
    \left[
    \CPEel{Q_1}{\alpha_1}
    \parallel
    \left(
        R_2
        \seriesconn
        \left(
            R_3
            \parallel
            \CPEel{Q_4}{\alpha_4}
        \right)
    \right)
    \right]
\]
}
\end{frame}

\begin{frame}{Simulation study - Battery data}
\begin{figure}
  \centering
  \includegraphics[width=0.75\linewidth]{/home/mitek/Work/Equation-Discovery-for-FODE/Results/campaigns/simulated_006_007_no_G_relaxed_5_relative_soft_l1_fscale_0p001_optimizer_value_lp_bic_1start/runs/simulated_simulated_spectrum_006/Results/Plots/Nyquists/Fits of simulated battery data.png}
  \caption{Simulated SOFC data fit results}    
\end{figure}

\end{frame}

\begin{frame}{Simulation study - SOFC data}
\begin{figure}
  \centering
  \includegraphics[width=0.75\linewidth]{/home/mitek/Work/Equation-Discovery-for-FODE/Results/campaigns/simulated_006_007_no_G_relaxed_5_relative_soft_l1_fscale_0p001_optimizer_value_lp_bic_1start/runs/simulated_simulated_spectrum_007/Results/Plots/Nyquists/Fits of simulated SOFC data.png}
  \caption{Simulated SOFC data fit results}    
\end{figure}

\end{frame}

\begin{frame}
  \frametitle{SOFC data}
\end{frame}

\begin{frame}
  \frametitle{Battery data}

\end{frame}

\begin{frame}
  \frametitle{TODO}
  \begin{itemize}
    \item Parallelization and fitting larger circuits
    \item Research \& implement alternative numerical algorithms
    \item Auto-tuning of the numerical algorithm parameters
    \item Experimental results over several battery cycles
    \item Compare it to some LLM type shit
  \end{itemize}
\end{frame}

\begin{frame}
  
\end{frame}

\appendix

\begin{frame}{Backup: parameters and residuals}
  For measured impedance $Z_i$ at frequency $f_i$,
  \[
    \widehat Z_i(\boldsymbol{\theta})
      = Z_{\mathrm{circuit}}(2\pi f_i;\boldsymbol{\theta}),
    \qquad
    e_i(\boldsymbol{\theta})
      = \widehat Z_i(\boldsymbol{\theta})-Z_i.
  \]
  \[
    r_{i,\mathrm{Re}}=\Re(e_i),
    \qquad
    r_{i,\mathrm{Im}}=\Im(e_i).
  \]
  The optimizer receives one real-valued residual vector:
  \[
    \widetilde{\mathbf r}(\boldsymbol{\theta})=
    \bigl(
      \widetilde r_{1,\mathrm{Re}},\ldots,\widetilde r_{N,\mathrm{Re}},
      \widetilde r_{1,\mathrm{Im}},\ldots,\widetilde r_{N,\mathrm{Im}}
    \bigr)^{\mathsf T}.
  \]
  Positive parameters use log space, while the CPE exponent stays linear:
  \[
    x_j=\log\theta_j,\quad \theta_j=e^{x_j},
    \qquad x_{\alpha}=\alpha.
  \]
\end{frame}

\begin{frame}{Backup: residual scaling}
  \small
  For component $c\in\{\mathrm{Re},\mathrm{Im}\}$,
  \[
    \widetilde r_{i,c}(\boldsymbol{\theta})
      = \frac{r_{i,c}(\boldsymbol{\theta})}{s_{i,c}}.
  \]
  \textbf{Unscaled and component-relative choices}
  \[
    s_{i,c}^{\mathrm{unscaled}}=1,
    \qquad
    s_{i,c}^{\mathrm{relative}}
      =\max\!\left(
        |[Z_i]_c|,
        0.1\,\operatorname{median}_{j}|[Z_j]_c|
      \right).
  \]
  \textbf{Repeated-measurement scatter choice}
  \[
    m_{i,c}=\operatorname{median}_{k}[Z_{ik}]_c,
    \qquad
    a_{i,c}=1.4826\,\operatorname{median}_{k}
      \left|[Z_{ik}]_c-m_{i,c}\right|,
  \]
  \[
    s_{i,c}^{\mathrm{scatter}}
      =\max\!\left(
        a_{i,c},
        \eta\,\operatorname{median}_{j}a_{j,c}
      \right).
  \]
\end{frame}

\begin{frame}{Backup: robust least squares}
  \small
  With $m=2N$ scalar residuals and $C=\texttt{f\_scale}$, SciPy's
  \texttt{soft\_l1} objective is
  \[
    F(\boldsymbol{\theta})
      =\frac{C^2}{2}\sum_{q=1}^{m}
       \rho\!\left(
         \frac{\widetilde r_q(\boldsymbol{\theta})^2}{C^2}
       \right),
    \qquad
    \rho(z)=2\left(\sqrt{1+z}-1\right).
  \]
  Our runs use
  \[
    C=\texttt{f\_scale}=10^{-3}.
  \]
  Residuals below $C$ behave approximately quadratically. Larger residuals
  receive progressively less weight.

  \vspace{0.5em}
  The reported first-round optimizer score is
  \[
    J_{\mathrm{opt}}=\frac{2F}{N}.
  \]
\end{frame}

\begin{frame}{Backup: fit and ranking scores}
  \small
  With $N$ frequencies and $k$ fitted circuit parameters,
  \[
    \mathrm{MSE}
      =\frac{1}{N}\sum_{i=1}^{N}
       \left(r_{i,\mathrm{Re}}^2+r_{i,\mathrm{Im}}^2\right),
  \]
  \[
    E_p
      =\frac{1}{N}\sum_{i=1}^{N}
       \left(|r_{i,\mathrm{Re}}|^p+|r_{i,\mathrm{Im}}|^p\right),
    \qquad p=0.5.
  \]
  \[
    \mathrm{BIC}=2N\log(\mathrm{MSE})+k\log(2N),
  \]
  \[
    L_{0.5}\text{-}\mathrm{BIC}
      =2N\log(E_{0.5})+k\log(2N).
  \]
  First-round fits enter refinement when
  \[
    J_i\le(1+0.20)J_{\min}=1.2J_{\min}.
  \]
\end{frame}

\end{document}
